\section{Geometric support, degeneracies, and domains}
\label{app:geometric-details}

We first examine how ties affect connector visibility and face dimension, then distinguish clearance at a connector point inside the face from distance to the face when that point lies outside it. The last two subsections explain how clipping and selected periodic images affect these checks.


\subsection{Weak visibility, persistent ties, and face dimension}
\label{app:geometry-degeneracy}

Consider a finite set of distinct sites in $\mathbb R^d$ with finite real weights. Fix $i\ne j$, and let $q=q_{ij}$ and $H=H_{ij}$ be the pairwise connector point and bisector defined in Section~4. Every point $x\in H$ can be written $x=q+u$ with $u\cdot\delta_{ij}=0$. As in Eq.~\eqref{eq:power-gap-displacement}, the competing-site inequalities restricted to $H$ are
\begin{equation}
2e_{ijk}\cdot u\le g_{ijk},
\qquad k\notin\{i,j\}.
\label{eq:appendix-face-halfspaces}
\end{equation}
Thus the common cell intersection is a (possibly empty or unbounded) closed convex polyhedron within $H$.

Weak connector visibility, $g_{ijk}\ge0$ for all $k$, is exactly the condition $q\in F_{ij}$; it need not imply that $q$ lies in the interior within the \emph{full} $(d-1)$-dimensional hyperplane $H$. Conversely, strict positivity of every competitor gap suffices for such interior membership, but is not necessary in the presence of persistent ties.

In particular, when $e_{ijk}=0$, inequality~\eqref{eq:appendix-face-halfspaces} is constant throughout $H$. If $g_{ijk}<0$, it excludes the entire bisector; if $g_{ijk}>0$, it imposes no restriction on the face; if $g_{ijk}=0$, it is an equality throughout $H$ and does not reduce the extent of $F_{ij}$ \emph{within} the bisector. In contrast, when $g_{ijk}=0$ and $e_{ijk}\ne0$, the corresponding halfspace boundary within $H$ passes through $q$. No ball of positive radius centered at $q$ within $H$ can then be contained in $F_{ij}$.

For an explicit persistent tie, take $p_i=(0,0)$, $p_j=(2,0)$, $p_k=(1,0)$ and $(w_i,w_j,w_k)=(0,0,-1)$. At every point of $H_{ij}=\{(1,y):y\in\mathbb R\}$, all three power distances agree. Consequently $F_{ij}=H_{ij}$ and $q_{ij}\in\operatorname{int}_{H_{ij}}F_{ij}$, yet $g_{ijk}=0$ identically. Here $C_i=\{(x,y):x\le1\}$ and $C_j=\{(x,y):x\ge1\}$ are full-dimensional cells sharing a facet, but the degenerate cell $C_k=H_{ij}$ also minimizes power there. Thus interior membership does not exclude a third minimizing site or imply strict visibility. If instead the selected pair is $i,k$, then $F_{ik}=H_{ik}=H_{ij}$ while $C_k$ has empty ambient interior; full-dimensionality of the intersection within its bisector does not by itself certify that both selected cells are full-dimensional.

We always use $\operatorname{int}_{H_{ij}}$, the interior relative to the \emph{entire pairwise bisector}, not the ordinary relative interior within $\operatorname{aff}(F_{ij})$. The latter is nonempty for every nonempty convex face, including a singleton, and would fail to distinguish full-dimensional interfaces.

\subsection{Centered clearance and exterior distance to a face}
\label{app:geometry-exterior}

The gap inequalities~\eqref{eq:appendix-face-halfspaces} also describe the distance to a realized face when the connector point itself is not feasible. Define
\begin{equation}
V_{ij}=\left\{u\in\delta_{ij}^{\perp}:
2e_{ijk}\cdot u\le g_{ijk}
\text{ for all }k\notin\{i,j\}\right\}.
\label{eq:face-displacement-feasible-set}
\end{equation}
Then
\begin{equation}
F_{ij}=q_{ij}+V_{ij}.
\label{eq:face-displacement}
\end{equation}
The set $V_{ij}$ is a closed convex polyhedron, possibly empty.

\paragraph{Interior clearance}
If $q_{ij}\in F_{ij}$, the largest radius of a closed Euclidean ball \emph{within} $H_{ij}$ centered at $q_{ij}$ and contained in $F_{ij}$ satisfies
\begin{equation}
c_{ij}=
\min_{\substack{k\notin\{i,j\}\\\|e_{ijk}\|>0}}
\frac{g_{ijk}}{2\|e_{ijk}\|},
\label{eq:weak-centered-clearance}
\end{equation}
with the minimum of an empty set interpreted as $+\infty$. Indeed, the minimum of each affine gap over the centered ball of radius $r$ is $g_{ijk}-2r\|e_{ijk}\|$; all these values must be nonnegative. The same argument works when $\sigma_{ij}=0$, not only when it is positive. A zero gap with nonzero $e_{ijk}$ forces $c_{ij}=0$; a zero gap with $e_{ijk}=0$ imposes no in-bisector restriction. When no constraint has a nonzero projected normal and $q_{ij}\in F_{ij}$, the whole bisector is contained in the face and $c_{ij}=+\infty$. Redundant halfspaces do not affect this derivation.

\paragraph{Exterior distance}
If $q_{ij}\notin F_{ij}$ but $F_{ij}\ne\varnothing$, define its \emph{exterior face distance} as the inherited Euclidean distance within the bisector,
\begin{equation}
d_{ij}^{\mathrm{out}}
=\operatorname{dist}_{H_{ij}}(q_{ij},F_{ij}).
\label{eq:exterior-face-distance}
\end{equation}
Using~\eqref{eq:face-displacement}, this distance is computed by the convex quadratic program
\begin{equation}
\boxed{
\begin{aligned}
(d_{ij}^{\mathrm{out}})^2
&=\min_{u\in\delta_{ij}^{\perp}}\|u\|^2,\\
\text{subject to}\quad
&2e_{ijk}\cdot u\le g_{ijk}
\qquad\text{for all }k\notin\{i,j\}.
\end{aligned}}
\label{eq:exterior-distance-qp}
\end{equation}
For a nonempty face the minimum is attained at a unique point: this is Euclidean projection onto a nonempty closed convex set. If the inequalities are infeasible, $F_{ij}=\varnothing$ and the standard distance-to-set convention gives $d_{ij}^{\mathrm{out}}=+\infty$; there is then no nearest point. The quantity measures distance to the \emph{specified pair's face}, not to the closest face of any site pair in the full diagram.

\paragraph{Why negative clearance ratios are not exterior distances}
Consider the formal algebraic extension
\begin{equation}
\widetilde c_{ij}=
\min_{\|e_{ijk}\|>0}\frac{g_{ijk}}{2\|e_{ijk}\|}.
\label{eq:formal-negative-clearance}
\end{equation}
For a violated inequality with nonzero projected normal, the positive quantity $-g_{ijk}/(2\|e_{ijk}\|)$ is the distance from $q_{ij}$ to the feasible \emph{halfspace for that single inequality}. The full face is the intersection of all halfspaces, so its distance from $q_{ij}$ can exceed each individual halfspace distance, and the intersection itself may be empty. Constant violated inequalities ($e_{ijk}=0,\ g_{ijk}<0$) are another obstruction: they make the whole face empty without appearing in~\eqref{eq:formal-negative-clearance}.

An explicit example with a nonempty face occurs in $\mathbb R^3$. Take
\begin{equation}
\begin{aligned}
p_1&=(0,0,0), & p_2&=(2,0,0),\\
p_3&=(1,-1,0), & p_4&=(1,0,-1),
\end{aligned}
\qquad
(w_1,w_2,w_3,w_4)=(0,0,2,2).
\label{eq:exterior-counterexample}
\end{equation}
Here $H_{12}=\{(1,y,z)\}$, $q_{12}=(1,0,0)$, and the two competitor inequalities require $y\ge1$ and $z\ge1$. The face is the nonempty quadrant $\{(1,y,z):y\ge1,\ z\ge1\}$, so
\begin{equation}
\widetilde c_{12}=-1,
\qquad d_{12}^{\mathrm{out}}=\sqrt2.
\label{eq:exterior-example-values}
\end{equation}
In the four-site configuration of Eq.~\eqref{eq:four-sites} with $(w_3,w_4)=(1,1)$, the face is instead empty despite a finite negative formal ratio. Hence a negative continuation of the centered-clearance formula cannot be identified with a signed distance to the true face. The centered interior radius and the distance to a realized face are therefore distinct geometric quantities.

\subsection{Clipping to a convex box}
\label{app:geometry-clipping}

For a convex box or, more generally, a fixed convex polytope
\begin{equation}
\Omega=\{x\in\mathbb R^d:a_\ell\cdot x\le b_\ell
\text{ for all }\ell\},
\label{eq:domain-halfspaces}
\end{equation}
the clipped intersection is $F_{ij}^{\Omega}=F_{ij}\cap\Omega$. Each domain inequality restricted to $x=q_{ij}+u\in H_{ij}$ takes the form
\begin{equation}
a_\ell\cdot u\le b_\ell-a_\ell\cdot q_{ij}.
\end{equation}
Let $P_{ij}$ be orthogonal projection onto $\delta_{ij}^{\perp}$. When $q_{ij}\in\Omega$, the largest radius of a centered in-bisector ball contained in $\Omega$ is
\begin{equation}
b_{ij}^{\Omega}=
\min_{\ell:\|P_{ij}a_\ell\|>0}
\frac{b_\ell-a_\ell\cdot q_{ij}}{\|P_{ij}a_\ell\|},
\label{eq:box-inplane-radius}
\end{equation}
with empty minimum $+\infty$; a constraint with zero projected normal has no effect on motion within $H_{ij}$ when already satisfied at $q_{ij}$.

If $q_{ij}\in F_{ij}\cap\Omega$, the centered radius of the clipped intersection is exactly $\min\{c_{ij},b_{ij}^{\Omega}\}$ by the two families of halfspace inequalities. The strict-visibility and interior-domain hypotheses in Section~4.4 guarantee that this value is positive; at the boundary, ties and domain-wall coincidences require explicit attention.

When a domain wall coincides with $H_{ij}$, a clipped intersection may have nonempty $(d-1)$-dimensional interior within $H_{ij}$ without forming an interface between two full-dimensional cells \emph{inside} $\Omega$. An interior-domain strict witness is a sufficient way to exclude this boundary-only behavior, but not a necessary condition for every useful clipped geometric intersection. No analogous finite halfspace formula is claimed here for arbitrary nonconvex clipping domains.

\subsection{Selected periodic images and competitor scope}
\label{app:geometry-periodic}

For a fixed lattice with translations $\Lambda t$, a geometric site instance has position $p_{k,t}=p_k+\Lambda t$ and retains base weight $w_k$. A periodic pair observation specifies two \emph{selected instances}, such as $(i,0)$ and $(j,s)$. The bisector and its connector point are computed for those instances; their competitors are \emph{all other instances}. In particular, other copies of base sites $i$ or $j$ must not be discarded merely because their base indices match one of the endpoints. Distinct base sites may be equivalent modulo the lattice: a coincident \emph{unselected} image produces a constant competitor power-gap inequality, which is redundant, excludes the selected cell, or creates a persistent tie depending on its sign. The observed selected endpoints must still be distinct as required in Section~2.

The same affine pointwise power inequalities and in-bisector projection arguments apply to each selected instance and competitor. However, an infinite collection of image inequalities is not automatically a finite polyhedron, and an arbitrary neighbor-image stencil does not establish a complete exact realization test or a uniform finite stability constant. An exact finite reduction requires a justified region of interest and relevant-image bound.

Let $V=\operatorname{span}\Lambda$ and let $P_V$ denote orthogonal projection onto $V$. Competition with a selected site's own images confines $P_V(x-p_i)$, for every $x$ in that image's power cell, to the compact lattice Voronoi region in $V$. For images of any competing base site, the $V^\perp$ component of squared distance is independent of the shift, and all images share the same weight. Thus, for each competing base site, its nearest image depends only on $P_Vx$, and finitely many shifts suffice over the compact projected region even when the nonperiodic directions are unbounded. This supplies a route to a finite relevant-image reduction for selected cells and their faces under fixed partial periodicity.

It does not justify an arbitrary fixed image stencil or a uniform competitor set for connector points that move without restriction as weights change. Changing lattice geometry also requires separate control. Section~4 therefore does not extend its finite-competitor stability estimate to an unrestricted periodic theorem.
